\documentclass[11pt]{article}

\usepackage[margin=1in]{geometry}
\usepackage{amsmath,amssymb,amsthm,mathtools}
\usepackage{booktabs}
\usepackage{microtype}
\usepackage{enumitem}
\usepackage{xcolor}
\usepackage[hidelinks]{hyperref}

\hypersetup{
  pdftitle={An Exact Locked-Permutation Certificate for Property B: m(5) >= 33},
  pdfauthor={Ayush Agarwal}
}

\newtheorem{theorem}{Theorem}
\newtheorem{lemma}[theorem]{Lemma}
\newtheorem{proposition}[theorem]{Proposition}
\newtheorem{corollary}[theorem]{Corollary}
\theoremstyle{remark}
\newtheorem{remark}[theorem]{Remark}

\newcommand{\Qlock}{\mathcal{Q}}
\newcommand{\ind}{\mathbf{1}}

\title{An Exact Locked-Permutation Certificate for Property B:\\
       \texorpdfstring{$m(5)\ge 33$}{m(5) >= 33}}
\author{Ayush Agarwal\\
\small Independent researcher, Bengaluru, India\\
\small \href{mailto:ayushag1002@gmail.com}{\texttt{ayushag1002@gmail.com}}\\
\small \href{https://orcid.org/0009-0000-7437-1929}%
{ORCID: 0009-0000-7437-1929}}
\date{}

\begin{document}
\maketitle

\begin{abstract}
Let $m(n)$ be the minimum number of edges in an $n$-uniform hypergraph
that is not two-colourable.  The previously established bounds at
uniformity five are $32\le m(5)\le51$.  We prove that every 5-uniform
hypergraph with at most 32 edges is two-colourable, and hence
\[
33\le m(5)\le51.
\]
The proof refines the random greedy method by locking a small,
deterministically selected set of low-degree vertices into prescribed
positions of a random permutation.  Pair compression, balanced-colouring
counting, and an unconditioned greedy estimate reduce a hypothetical
obstruction to five possible vertex counts, 21 through 25.  Degree and
codegree selection lemmas then reduce those cases to finite trace-profile
calculations on at most five selected vertices.  The exact computations
include 3,081 certificate records for vertex counts 21 and 22, together with
six labelled clique partitions of $K_4$ and 32 of $K_5$.  A separately
written implementation regenerates the finite universes and re-derives the
probability coefficients.  We also determine a precise ceiling for the
current certificate families: exhaustive optimization of their locked
positions does not extend this argument to $m(5)\ge34$.
\end{abstract}

\noindent\textbf{2020 Mathematics Subject Classification.}
05C15 (primary); 05C65, 05D40 (secondary).

\smallskip
\noindent\textbf{Keywords.}
Property B; hypergraph two-colouring; random greedy colouring;
computer-assisted proof; exact rational certificate; set splitting.

\section{Introduction}\label{sec:introduction}

An $n$-uniform hypergraph has \emph{Property B} if its vertices admit a
red--blue colouring with no monochromatic edge.  For integers $n$ and $v$,
write
\[
m(n,v):=\min\{|E(H)|:H\text{ is an $n$-uniform non-two-colourable
hypergraph on $v$ vertices}\},
\]
with $m(n,v)=\infty$ if no such hypergraph exists, and put
$m(n)=\min_v m(n,v)$.  Equivalently, Property B is the feasibility version
of monotone NAE-SAT or set splitting with every constraint of size $n$.

The exact values of $m(n)$ are known only for $n\le4$
\cite{GrillLinzmayerOverview2026}; in particular,
$m(4)=23$ was established by an exhaustive computation of
\"Osterg{\aa}rd \cite{Ostergard2014}.  At uniformity five, Abbott and Hanson
constructed a 51-edge obstruction \cite{AbbottHanson1969}.  On the lower
side, Aglave, Amarnath, Shannigrahi and Singh proved $m(5)\ge29$
\cite{AglaveEtAl2020}, and Grill and Linzmayer subsequently proved
$m(5)\ge32$ \cite{GrillLinzmayer2024}.  Thus the interval immediately
preceding this work was
\[
32\le m(5)\le51.
\]
The 2026 overview of Grill and Linzmayer cites that 2024 lower-bound preprint
\cite{GrillLinzmayerOverview2026}.  To the best of our knowledge, as
of August 23, 2026 no publicly available proof of $m(5)\ge33$ had appeared.
A dated primary-source search is included in the supplementary artifact; as
usual, such a search cannot rule out unpublished work.

Our main result improves the lower endpoint by one.

\begin{theorem}\label{thm:main}
Every 5-uniform hypergraph with at most $32$ edges has Property B.
Consequently,
\[
m(5)\ge33.
\]
\end{theorem}

The proof belongs to the random greedy line initiated by Pluh\'ar
\cite{Pluhar2009}, developed asymptotically by Cherkashin and Kozik
\cite{CherkashinKozik2015}, and adapted to exact small parameters by Grill
and Linzmayer \cite{GrillLinzmayer2024}.  The comparison with the last of
these is precise.  Their unconditioned discrete estimate and low-degree
conditioning prove two-colourability for hypergraphs with at most 31 edges.
Here we give a self-contained
normalization for 32 edges, derive a general formula for locking any fixed
number of vertices, and combine three-, four- and five-lock certificates
with structural selection arguments.  The earlier $m(5)\ge32$ computation
is not used as a premise.

The computer-assisted portion of the proof of
Theorem~\ref{thm:main} is deliberately smaller than an enumeration of
hypergraphs.  It performs the following tasks only:
\begin{enumerate}[leftmargin=2em]
\item enumerate a safe superset of the integer trace profiles of three
      selected vertices;
\item enumerate clique partitions of $K_4$ and $K_5$ that encode traces on
      pairwise-codegree-one selected sets;
\item evaluate a displayed rational expression at finite, explicitly listed
      position assignments; and
\item compare each reduced fraction with one.
\end{enumerate}
No SAT solver, mixed-integer solver, stochastic search, or floating-point
acceptance test enters Theorem~\ref{thm:main}.  Proposition~\ref{prop:finite}
states the complete machine-checked interface; Sections~\ref{sec:cases}
and~\ref{sec:artifact} explain how it is independently replayed.

\section{A pair-covered 32-edge normal form}\label{sec:normal}

We work with finite simple hypergraphs.  Duplicate edges may always be
discarded.  The following standard pair-identification reduction appears in
the exact Property-B literature, including Goldberg and Russell
\cite{GoldbergRussell1995}; we include the proof because the repeated padding
step makes the present theorem independent of prior numerical bounds.

\begin{lemma}[pair compression and padding]\label{lem:compression}
If a non-two-colourable 5-uniform hypergraph with at most $32$ edges exists,
then there is a non-two-colourable simple 5-uniform hypergraph with exactly
$32$ edges in which every pair of vertices occurs in an edge.
\end{lemma}

\begin{proof}
A 5-uniform hypergraph on at most eight vertices is two-colourable: split
its vertex set into two classes of size at most four.  Hence an obstruction
has at least nine vertices.  Since $\binom95=126>32$, distinct edges can be
added until there are exactly 32; adding constraints preserves
non-two-colourability.

Suppose that vertices $x,y$ occur together in no edge.  Identify $x$ and
$y$.  No image edge loses a vertex, so every image edge still has size five.
Any two-colouring of the quotient pulls back by giving $x$ and $y$ the
quotient colour, so the quotient remains non-two-colourable.  After
coincident image edges are discarded it has at most 32 edges.  If its order
is at most eight, the preceding split is a contradiction; otherwise it has
at least nine vertices and can again be padded to 32 distinct edges.
Repeating this operation strictly decreases the order at every
identification, so it terminates at the claimed pair-covered hypergraph.
\end{proof}

Let $H$ henceforth be such a normal-form hypergraph, and write
\[
v=|V(H)|,\qquad d(x)=|\{e:x\in e\}|,
\qquad \lambda_{xy}=|\{e:\{x,y\}\subseteq e\}|.
\]
Pair coverage gives
\begin{equation}\label{eq:pair-cover}
\lambda_{xy}\ge1,
\qquad
\binom v2\le32\binom52=320,
\end{equation}
so $v\le25$.

A uniformly chosen colouring with colour-class sizes
$\lfloor v/2\rfloor$ and $\lceil v/2\rceil$ makes any fixed 5-edge
monochromatic with probability
\[
\frac{\binom{\lfloor v/2\rfloor}5+\binom{\lceil v/2\rceil}5}
     {\binom v5}.
\]
The first-moment method therefore gives the familiar bound
\begin{equation}\label{eq:balanced}
m(5,v)\ge
\frac{\binom v5}
{\binom{\lfloor v/2\rfloor}5+\binom{\lceil v/2\rceil}5};
\end{equation}
see also \cite{GoldbergRussell1995,GrillLinzmayer2024}.  For
$10\le v\le18$, the ceilings of the right-hand side are
\[
126,66,66,48,48,39,39,34,34,
\]
respectively; orders at most nine are immediate from the same balanced
split.  Hence
\begin{equation}\label{eq:range-initial}
19\le v\le25.
\end{equation}

We shall repeatedly use
\begin{equation}\label{eq:degree}
d(x)\ge\left\lceil\frac{v-1}{4}\right\rceil,
\qquad
\sum_xd(x)=5|E(H)|=160,
\end{equation}
and the exact local pair excess
\begin{equation}\label{eq:rho}
\rho_x:=\sum_{y\ne x}(\lambda_{xy}-1)
=4d(x)-(v-1).
\end{equation}
The first inequality in \eqref{eq:degree} follows because every edge through
$x$ covers four of the $v-1$ pairs incident with $x$.

\section{The locked-permutation bound}\label{sec:locked}

Choose an ordered set $S=(s_1,\ldots,s_s)$ of \emph{locked} vertices, where
$0\le s\le5$ and $f:=v-s\ge9$, and distinct positions
$p_1,\ldots,p_s\in[v]$.  Assign the
other vertices uniformly to the free positions.  For $A\subseteq[s]$, let
\[
l_A=|\{e\in E(H):e\cap S=\{s_i:i\in A\}\}|.
\]
For a position $k$, let $b(k)$ and $a(k)$ be the numbers of free positions
before and after $k$.  Write $A\prec k$ if every lock in $A$ precedes $k$,
and $A\succ k$ if every lock in $A$ follows $k$.  Empty conditions are
vacuous, and binomial coefficients outside their natural range are zero.

Run Pluh\'ar's greedy algorithm: begin with every vertex red, process the
permutation from first to last, and turn the current vertex blue precisely
when it is last in an otherwise-red edge.  No red edge survives.  If a blue
edge survives, its earliest vertex is simultaneously last in a red witness
edge and first in the blue edge.  These two edges meet only at that vertex;
call it \emph{critical}.

For a free position $k$, union bounds for being last, first, or both are
\begin{align}
L_k&=\sum_{\substack{A\prec k\\|A|\le 4}}l_A\frac{5-|A|}{f}
 \frac{\binom{b(k)}{4-|A|}}{\binom{f-1}{4-|A|}},\label{eq:free-last}\\
F_k&=\sum_{\substack{A\succ k\\|A|\le 4}}l_A\frac{5-|A|}{f}
 \frac{\binom{a(k)}{4-|A|}}{\binom{f-1}{4-|A|}},\label{eq:free-first}\\
B_k&=\sum_{\substack{A\cap B=\varnothing\\A\prec k,\ B\succ k\\
|A|,|B|\le 4}}
\left(l_Al_B-\ind_{A=B}l_A\right)
\frac{\binom{b(k)}{4-|A|}\binom{a(k)}{4-|B|}}
{f\binom{f-1}{4-|A|}\binom{f-1-(4-|A|)}{4-|B|}}.
\label{eq:free-both}
\end{align}
Trace classes with $|A|=5$ contain no free vertex and therefore do not
contribute at a free position.
For a locked position $k=p_i$, the corresponding expressions are
\begin{align}
L_k&=\sum_{\substack{i\in A\\A\setminus\{i\}\prec k}}
l_A\frac{\binom{b(k)}{5-|A|}}{\binom f{5-|A|}},
\label{eq:locked-last}\\
F_k&=\sum_{\substack{i\in A\\A\setminus\{i\}\succ k}}
l_A\frac{\binom{a(k)}{5-|A|}}{\binom f{5-|A|}},
\label{eq:locked-first}\\
B_k&=\sum_{\substack{A\cap B=\{i\}\\
A\setminus\{i\}\prec k,\ B\setminus\{i\}\succ k}}
\left(l_Al_B-\ind_{A=B}l_A\right)
\frac{\binom{b(k)}{5-|A|}\binom{a(k)}{5-|B|}}
{\binom f{5-|A|}\binom{f-(5-|A|)}{5-|B|}}.
\label{eq:locked-both}
\end{align}

\begin{lemma}[locked-permutation lemma]\label{lem:locked}
With the notation above,
\begin{equation}\label{eq:Q}
\Pr(\text{greedy failure})\le
\Qlock(H,S,p):=\sum_{k=1}^{v}\min\{L_k,F_k,B_k\}.
\end{equation}
In particular, $\Qlock(H,S,p)<1$ implies that $H$ has Property B.
\end{lemma}

\begin{proof}
For an edge of type $A$ at a free position, the factor $(5-|A|)/f$
chooses its occupant and the binomial ratio places its remaining free
vertices on the required side; this proves \eqref{eq:free-last} and
\eqref{eq:free-first}.  A compatible ordered witness pair must have disjoint
locked traces.  After its common free vertex occupies $k$, the other free
vertices of the two edges occupy disjoint earlier and later slots, giving
\eqref{eq:free-both}.  The category product is an upper bound because it may
also count edge pairs having an additional common free vertex.  The diagonal
subtraction only forbids using one edge twice.

At a locked position the common vertex is fixed, yielding
\eqref{eq:locked-last}--\eqref{eq:locked-both}; now the two locked traces
must meet exactly in the locked vertex.  Criticality at $k$ is contained in
each of the last-edge, first-edge and paired events, so its probability is
at most their minimum.  Summing over positions proves \eqref{eq:Q}.
\end{proof}

The locked set and its position assignment may depend deterministically on
the fixed hypergraph and its trace profile.  Only after this choice do we
randomize the remaining vertices.  Consequently, different profiles may
use different assignments without any additional union bound.

The zero-lock specialization of Lemma~\ref{lem:locked} is useful at the two
smallest remaining orders.  Let $\gamma$ be the number of \emph{ordered}
pairs of distinct edges meeting in exactly one vertex.  Then
\begin{equation}\label{eq:zero-lock}
U_0(v,\gamma)=\sum_{k=1}^{v}\frac1v\min\!\left\{
160\frac{\binom{k-1}4}{\binom{v-1}4},
160\frac{\binom{v-k}4}{\binom{v-1}4},
\gamma\frac{\binom{k-1}4\binom{v-k}4}
{\binom{v-1}4\binom{v-5}4}
\right\}.
\end{equation}
Using $\gamma\le32\cdot31=992$ gives
\[
U_0(19,992)=\frac{222656}{230945}<1.
\]
For $v=20$, the pair-incidence surplus is
\[
R=320-\binom{20}{2}=130.
\]
Moreover,
\[
\sum_{x<y}\binom{\lambda_{xy}}2
=\sum_{\{e,f\}}\binom{|e\cap f|}{2}\ge R.
\]
Distinct 5-edges meet in at most four vertices.  Thus at least
$\lceil130/6\rceil=22$ unordered edge pairs meet in at least two vertices,
so $\gamma\le992-2\cdot22=948$, and
\[
U_0(20,948)=\frac{20874}{20995}<1.
\]
It remains to exclude
\begin{equation}\label{eq:remaining-v}
21\le v\le25.
\end{equation}

\section{Finite trace profiles and certification}\label{sec:profiles}

This section specifies the finite computation in enough detail for an
independent implementation.

\subsection{Three-lock profiles}

For three selected vertices, counts are stored in binary-mask order
\begin{equation}\label{eq:mask-order}
(l_{000},l_{001},l_{010},l_{011},l_{100},l_{101},l_{110},l_{111}),
\end{equation}
where the $i$th selected vertex is represented by bit $2^{i-1}$.  Thus, for
example, $l_{011}$ counts edges containing the first two selected vertices
and not the third.

\begin{lemma}[complete safe profile universe]\label{lem:profiles}
Fix selected degrees $(d_1,d_2,d_3)$.  Every trace profile of a
pair-covered 32-edge hypergraph is produced by the following enumeration:
choose nonnegative integers
\[
t=l_{111},\quad q_{12}=l_{011},\quad
q_{13}=l_{101},\quad q_{23}=l_{110},
\]
derive
\begin{align*}
l_{001}&=d_1-q_{12}-q_{13}-t,\\
l_{010}&=d_2-q_{12}-q_{23}-t,\\
l_{100}&=d_3-q_{13}-q_{23}-t,\\
l_{000}&=32-\sum_{\alpha\ne000}l_\alpha,
\end{align*}
and retain the tuple precisely when every entry is nonnegative and
\[
q_{12}+t\ge1,\qquad q_{13}+t\ge1,
\qquad q_{23}+t\ge1.
\]
The output may contain unrealizable profiles, but cannot omit a realizable
one.
\end{lemma}

\begin{proof}
For an actual profile, $t,q_{12},q_{13},q_{23}$ are among the loop choices.
The three displayed degree equations uniquely force the singleton counts,
the 32-edge total uniquely forces $l_{000}$, and pair coverage gives the
three final inequalities.  Hence every actual profile is retained.  Since
no constraints involving unselected vertices are imposed, the output is in
general a safe superset rather than a realizability characterization.
\end{proof}

If $d_1\le d_2\le d_3$ are the three smallest degrees, then
\begin{equation}\label{eq:degree-triples}
\begin{aligned}
d_1&\ge\left\lceil\frac{v-1}{4}\right\rceil,\\
d_1+d_2+(v-2)d_3&\le160.
\end{aligned}
\end{equation}
Looping over all ordered integer triples satisfying \eqref{eq:degree-triples}
therefore gives the complete degree universe used below.

\subsection{Clique partitions}

\begin{lemma}[trace partitions]\label{lem:trace}
If every pair in a selected set $S$ has codegree one, then the traces
$e\cap S$ of size at least two partition $E(K_S)$ into edge sets of cliques.
Once this clique partition and the selected degrees are fixed, all singleton
and empty trace counts are forced.
\end{lemma}

\begin{proof}
A trace $T$ contributes precisely the pairs in $\binom T2$.  Every selected
pair occurs in one trace, so these clique edge sets are disjoint and cover
$E(K_S)$.  Subtracting the incidences of the non-singleton traces from each
selected degree determines the singleton counts, and the total edge count
then determines the empty count.
\end{proof}

The clique-partition generator is complete by construction: choose the
lexicographically first uncovered pair, branch over every clique containing
that pair whose edge set is disjoint from the already covered pairs, and
recurse.  It returns six labelled partitions of $K_4$ and 32 labelled
partitions of $K_5$.  The $K_4$ partitions have the three types
\[
K_4,\qquad K_3+3K_2,\qquad 6K_2.
\]

\subsection{Pseudocode and position menus}

The primary computation is summarized below.  All arithmetic in
\textsc{LockedBound} is integer arithmetic and reduced rational arithmetic
using equations \eqref{eq:free-last}--\eqref{eq:Q}.

\begin{verbatim}
PROFILE3(d1,d2,d3):
  for t,q12,q13,q23 in their finite degree-bounded ranges:
    derive the three singleton counts from d1,d2,d3
    derive the empty count from the 32-edge total
    if all counts are nonnegative and each selected pair is covered:
      emit the eight-count tuple

CERTIFY(v, profile, menu):
  best := +infinity
  for injective labelled position assignment p in menu:
    best := min(best, LOCKEDBOUND(v,profile,p))
  return best and the lexicographically first minimizing p

CLIQUEPARTITIONS(s, covered):
  if every pair of K_s is covered: emit the selected cliques
  otherwise let xy be the first uncovered pair
  for each clique C containing xy and no covered pair:
    CLIQUEPARTITIONS(s, covered union E(C))
\end{verbatim}

For a tuple $(a_1,\ldots,a_s)$, ``all permutations'' below means all
label-to-position bijections onto that position set.  The complete menus are
\begin{center}
\begin{tabular}{@{}lll@{}}
\toprule
order & profile family & allowed assignments\\
\midrule
21 & three lowest & all permutations of
  $(10,11,12),(1,2,3),$\\
   & & $(11,12,13),(1,2,12)$\\
22 & three lowest/Ramsey & all permutations of
  $(10,11,12),(1,2,3),(1,2,11),$\\
   & & $(1,2,13),(1,11,13)$\\
23 & three lowest & all permutations of $(11,12,13)$\\
23 & degrees $(6,6,7,7)$ & all permutations of $(10,11,12,13)$ or
  $(12,13,14,15)$\\
24 & degrees $(6,6,6,6)$ & all permutations of $(11,12,13,14)$\\
25 & degrees $(6,6,6,6,6)$ & all permutations of $(11,12,13,14,15)$\\
\bottomrule
\end{tabular}
\end{center}
For the five-lock $v=23$ family with degrees $(6,7,7,7,7)$, the complete
menu is the following list of eleven \emph{ordered} assignments:
\begin{equation}\label{eq:menu-v23}
\begin{gathered}
(13,16,17,18,14),\ (11,4,3,2,1),\ (16,15,17,18,14),\\
(10,4,3,1,2),\ (15,16,17,18,14),\ (10,4,3,2,1),\\
(12,15,16,17,14),\ (17,15,16,18,14),\\
(12,16,17,18,14),\ (10,15,16,17,14),\ (12,4,3,2,1).
\end{gathered}
\end{equation}
No unlisted position is used in the proof.
Among the six $v=21$ position bases screened during optimization, an
exhaustive exact check of all base subsets shows that no three-base subset
attains the worst-case value $9599/9724$; the four-base menu displayed above
does.  This is a minimality statement within that candidate family, not
among all conceivable three-lock position sets.

\begin{proposition}[exact finite certificate]\label{prop:finite}
The algorithms above, evaluated with exact rational arithmetic, give the
following exhaustive results.
\begin{center}
\begin{tabular}{@{}llrrl@{}}
\toprule
$v$ & family & profiles & residuals & largest closing bound\\
\midrule
21 & three lowest & 1,864 & 0 & $9599/9724$\\
22 & three lowest & 1,000 & 3 & ---\\
22 & Ramsey independent & 2 & 0 & $91231/92378$\\
22 & Ramsey triangle & 215 & 0 & $456463/461890$\\
23 & degrees $(6,6,6)$ & 192 & 0 & $245/247$\\
23 & degrees $(6,6,7)$ & 220 & 2 & ---\\
23 & degrees $(6,7,7)$ & 263 & 17 & ---\\
23 & admissible $K_4$ partitions & 3 & 0 & $8315/8398$\\
23 & all $K_5$ partitions & 32 & 0 & $45784/51051$\\
24 & all $K_4$ partitions & 6 & 0 & $16671/16796$\\
25 & favourable $K_5$ partitions & 7 & 0 & $4188/4199$\\
\bottomrule
\end{tabular}
\end{center}
A dash denotes a family containing residual profiles, for which no
family-wide strict closing bound exists.
Here ``residual'' means that no assignment in the stated menu gives a value
strictly below one.  In the three $v=22$ residuals the selected degrees are
$(7,7,7)$, the pair-codegree multiset is $(1,1,2)$, and $l_{111}=0$.  In the
$v=23$ residuals, the two $(6,6,7)$ profiles have all selected codegrees one,
and all 17 $(6,7,7)$ profiles have selected codegrees at most two.
\end{proposition}

The proposition is a direct output assertion of the primary generator and
is independently recomputed by a second implementation.  To make the scale
and encoding concrete, three representative closing rows are
\begin{center}
\begin{tabular}{@{}llll@{}}
\toprule
family & degrees & counts in \eqref{eq:mask-order}; positions & bound\\
\midrule
$v=21$ worst & $(7,7,7)$ &
$(15,5,4,1,4,1,2,0);(1,2,12)$ & $9599/9724$\\
$v=22$ independent & $(7,7,7)$ &
$(14,5,5,1,5,1,1,0);(1,2,3)$ & $91231/92378$\\
$v=22$ triangle & $(7,7,7)$ &
$(17,3,3,2,3,2,2,0);(1,11,13)$ & $456463/461890$\\
\bottomrule
\end{tabular}
\end{center}
The complete records, not merely these representatives, are included in the
artifact described in Section~\ref{sec:artifact}.

\section{The five remaining orders}\label{sec:cases}

We now combine Proposition~\ref{prop:finite} with deterministic selection
arguments.  This separation is important: the computer checks profiles and
fractions, while the proofs below show that every hypothetical hypergraph
contains a selected set belonging to one of the checked families.

\subsection{Order 21}

Let $d_1\le d_2\le d_3$ be the three smallest degrees.  From
\eqref{eq:degree-triples},
\[
d_1\ge5,\qquad d_1+d_2+19d_3\le160.
\]
The complete list is
\[
\begin{split}
&(5,5,5),(5,5,6),(5,5,7),(5,6,6),(5,6,7),\\
&(5,7,7),(6,6,6),(6,6,7),(6,7,7),(7,7,7).
\end{split}
\]
Lemma~\ref{lem:profiles} gives 1,864 profiles, all closed by the first line
of Proposition~\ref{prop:finite}.  Therefore $v=21$ is impossible.

\subsection{Order 22}\label{subsec:v22}

The complete lowest-degree triple list is
\[
(6,6,6),(6,6,7),(6,7,7),(7,7,7),
\]
which produces 1,000 safe profiles.  The menu closes 997.  The three
residuals are exactly the labelings
\begin{equation}\label{eq:v22-residuals}
\begin{split}
&(15,4,4,2,5,1,1,0),\\
&(15,4,5,1,4,2,1,0),\\
&(15,5,4,1,4,1,2,0).
\end{split}
\end{equation}
Their best menu value is $1395847/1385670>1$; they are not claimed to be
realisable.  What matters is the common structural consequence: the selected
vertices have degrees $(7,7,7)$, pair-codegrees $(1,1,2)$ in some order, and
no edge contains all three.

Because the selected vertices were the three smallest, a residual
obstruction has minimum degree seven.  The total excess above seven is
\[
160-22\cdot7=6,
\]
so at most six vertices have degree greater than seven and at least 16 have
degree exactly seven.  On the degree-seven vertices define
\[
xy\in E(G)\quad\Longleftrightarrow\quad\lambda_{xy}\ge2.
\]
Because the minimum degree is seven, any three degree-seven vertices may be
chosen as the three lowest-degree vertices in the preceding enumeration.
Every graph on at least six vertices contains either an independent triple
or a triangle, since $R(3,3)=6$.  In the independent case, pair coverage
forces all three selected codegrees to equal one; in the triangle case all
three are at least two.  Lemma~\ref{lem:profiles} produces exactly two safe
profiles in the first class and 215 in the second, and Proposition
\ref{prop:finite} closes both classes strictly.  Hence $v=22$ is impossible.

\subsection{Order 23}\label{subsec:v23}

The degree constraints leave only
\[
(d_1,d_2,d_3)=(6,6,6),(6,6,7),(6,7,7).
\]
The central three-lock calculation has profile/residual counts
$192/0$, $220/2$ and $263/17$, respectively.  Thus a residual obstruction
has no three degree-six vertices.  It must nevertheless have a degree-six
vertex, since $23\cdot7=161>160$.  If there are $a$ degree-six vertices,
then
\[
\sum_x(d(x)-7)=-1.
\]
With $1\le a\le2$, the only possible degree sequences are therefore
\begin{equation}\label{eq:v23-sequences}
(6^2,7^{20},8),\qquad(6,7^{22}).
\end{equation}
Ties among degree-seven vertices may be broken arbitrarily, so the residual
conclusions apply to every degree-seven vertex or pair chosen below.

First consider $(6^2,7^{20},8)$.  Let $u,v$ be the degree-six vertices.
For every degree-seven vertex $y$, the selected triple $(u,v,y)$ must be one
of the two $(6,6,7)$ residual profiles; otherwise it is already closed.
Proposition~\ref{prop:finite} therefore gives
\[
\lambda_{uv}=\lambda_{uy}=\lambda_{vy}=1
\quad\text{for every such }y.
\]
Let $e$ be the unique edge containing $u,v$.  At most three of the 20
degree-seven vertices lie in $e$, so at least 17 lie outside it.  Among those
17, some pair $y,z$ has $\lambda_{yz}=1$: if every such pair were repeated,
then each vertex in this set would have local excess at least 16, contradicting
\[
\rho_y=4\cdot7-22=6.
\]
Thus $S=\{u,v,y,z\}$ is pairwise codegree one, and the unique trace
containing $u,v$ is exactly $\{u,v\}$ because $y,z\notin e$.  Of the six
labelled $K_4$ clique partitions, precisely three have this property.  The
four-lock row of Proposition~\ref{prop:finite} closes all three.  The
$6K_2$ row uses positions $(12,13,14,15)$ and has value $8315/8398$; the two
$K_3+3K_2$ rows have value $4152/4199$.  All are strict.

Now consider $(6,7^{22})$, and let $u$ be the degree-six vertex.  Applying
the $(6,7,7)$ residual consequence to $u$ and every pair of degree-seven
vertices shows that every pair-codegree in $H$ is at most two.  Pair coverage
therefore makes every codegree one or two.  Define the repeated-pair graph
\[
xy\in E(G)\quad\Longleftrightarrow\quad\lambda_{xy}=2.
\]
Equation \eqref{eq:rho} now equals the graph degree, so $G$ has degree
sequence $(2,6^{22})$ and 67 edges.  Let $A=N_G(u)$ and
$B=V(H)\setminus(\{u\}\cup A)$.  Then $|A|=2$, $|B|=20$, and for
$R=\{u\}\cup A$ we have $e(G[R])\in\{2,3\}$.  The number of graph edges
incident with $R$ is
\[
\sum_{x\in R}d_G(x)-e(G[R])=14-e(G[R]),
\]
whence
\[
e(G[B])=67-(14-e(G[R]))\in\{55,56\}.
\]
The Caro--Wei inequality \cite{Caro1979,Wei1981}, followed by
Cauchy--Schwarz, gives
\[
\alpha(G[B])\ge\sum_{x\in B}\frac1{d_B(x)+1}
\ge\frac{|B|^2}{2e(G[B])+|B|}
\ge\frac{400}{132}=\frac{100}{33}>3.
\]
Choose four independent vertices in $B$.  Together with $u$ they form a
five-set with every pair-codegree one.  Lemma~\ref{lem:trace} gives one of
the 32 labelled $K_5$ partitions, and the explicit menu
\eqref{eq:menu-v23} closes all 32 with maximum $45784/51051<1$.  This
excludes both sequences in \eqref{eq:v23-sequences}.

\subsection{Order 24}

Every degree is at least six, and the total excess above six is
\[
160-24\cdot6=16.
\]
Hence at least eight vertices have degree six; call their set $D$.  For
$x\in D$, equation \eqref{eq:rho} gives $\rho_x=1$.  The graph induced on
$D$ by repeated pairs consequently has maximum degree one, and therefore
has an independent set of size at least $\lceil|D|/2\rceil\ge4$.  The four
selected vertices are pairwise codegree one, so their trace is one of the
six labelled $K_4$ clique partitions.  Proposition~\ref{prop:finite} closes
all six with maximum $16671/16796<1$.

\subsection{Order 25}\label{subsec:v25}

The degree excess above six is now
\[
160-25\cdot6=10,
\]
so the degree-six set $D$ has size at least 15.  For $x\in D$,
\eqref{eq:rho} gives $\rho_x=0$; every pair incident with $D$ therefore has
codegree exactly one.

Suppose first that an edge $e$ contains at least four points of $D$.  Choose
four of them and a fifth vertex $z\in D\setminus e$, which is possible since
$|D|\ge15$.  The trace of $e$ on this selected set is a $K_4$.  Each of the
four pairs joining $z$ to that $K_4$ must be a separate $K_2$ trace, because
combining two would reuse their internal $K_4$ pair.  Thus the trace type is
exactly $K_4+4K_2$.

Otherwise every edge contains at most three points of $D$, and hence covers
at most one triple of $D$.  There are at most 32 covered triples.  If every
five-set of $D$ contained a covered triple, double-counting pairs $(T,S)$
with $T\subseteq S$, $|T|=3$, $|S|=5$, would give
\[
\binom{|D|}{5}\le32\binom{|D|-3}{2}.
\]
But
\[
\frac{\binom{|D|}{5}}{\binom{|D|-3}{2}}
=\frac{\binom{|D|}{3}}{10}
\ge\frac{\binom{15}{3}}{10}=\frac{91}{2}>32,
\]
a contradiction.  Therefore some five-set $S\subseteq D$ contains no
covered triple.  Since all its pairs have codegree one, every nonsingleton
trace has size two and the trace type is exactly $10K_2$.

The favourable-profile computation contains one $10K_2$ profile, five
labelled $K_4+4K_2$ profiles, and one $K_5$ regression profile.  The
selection above needs only the first two types; all are closed, and the
largest value among the seven is $4188/4199<1$.  This excludes $v=25$ and
completes the proof of Theorem~\ref{thm:main}.

\section{A scoped optimization ceiling}\label{sec:ceiling}

It is natural to ask whether a richer search over locked positions, with no
new combinatorial idea, would prove $m(5)\ge34$.  The following exhaustive
calculation answers that narrower question negatively.  Write
$\Qlock_{33}(l,p)$ for the evaluator in \eqref{eq:Q} applied to a trace
profile $l$ whose counts sum to 33.  The result below is not an upper bound
on $m(5)$ and does not assert that the displayed safe profiles are
realisable.

For context, the same normalization with 33 edges leaves exactly the orders
considered below.  Pair coverage gives $\binom v2\le330$, hence $v\le26$,
and \eqref{eq:balanced} excludes $v\le18$.  At $v=19$, pair-incidence
surplus gives at least $\lceil(330-\binom{19}{2})/6\rceil=27$ unordered edge
pairs meeting in at least two vertices.  Thus $\gamma\le33\cdot32-2\cdot
27=1002$.  Substituting $5\cdot33=165$ for $160$ in the right-hand side of
\eqref{eq:zero-lock}, the corresponding zero-lock failure-probability bound is
\[
\frac{228291}{230945}<1.
\]
At $v=26$, pair coverage forces minimum degree seven, contradicting
$26\cdot7>5\cdot33=165$.  A 33-edge extension of the method must therefore
handle precisely $20\le v\le25$.

\begin{proposition}[ceiling of position optimization]\label{prop:ceiling}
Consider the locked-permutation bound \eqref{eq:Q} using only selected
degrees and their trace counts.
\begin{enumerate}[label=(\roman*),leftmargin=2em]
\item For a hypothetical 33-edge core on 25 vertices, the safe five-lock
profile with selected degrees $(6,6,6,6,6)$ and trace partition $10K_2$
satisfies
\[
\min_{\substack{p_1,\ldots,p_5\in[25]\\\text{distinct}}}
\Qlock_{33}(l,p)=\frac{8649}{8398}>1.
\]
This minimum was obtained over all
$25\cdot24\cdot23\cdot22\cdot21=6{,}375{,}600$ labelled placements; symmetry
reduces the enumeration to $\binom{25}{5}=53{,}130$ position sets.
For comparison, the 32-edge certificate used in
Section~\ref{subsec:v25} gives $4188/4199<1$ for the corresponding profile.
\item Each of the five labelled $K_4+4K_2$ profiles forced by the other
branch of the selection in Section~\ref{subsec:v25} has global minimum
\[
\min_p\Qlock_{33}(l,p)=\frac{325}{323}>1.
\]
One representative minimizing assignment is $(12,1,13,14,15)$.  The
automorphism group has order 24, so 265,650 quotient representatives cover
all 6,375,600 labelled placements for each labelled trace profile.  At 32
edges the corresponding closing value is $37/38<1$.
\item At each order $v=20,21,\ldots,25$, the safe three-lock universe for a
33-edge core contains a profile whose global minimum over all labelled
three-lock placements is at least one.
\end{enumerate}
The exact witnesses for the last assertion are listed in mask order
\eqref{eq:mask-order}:
\begin{center}
\small
\begin{tabular}{@{}clll@{}}
\toprule
$v$ & degrees; counts & minimizing positions & global minimum\\
\midrule
20 & $(8,8,8);(15,4,4,2,4,2,2,0)$ & $(1,2,12)$ & $2287/2210$\\
21 & $(7,7,7);(15,5,5,1,5,1,1,0)$ & $(1,2,3)$ & $119365/116688$\\
22 & $(7,7,7);(15,5,5,1,5,1,1,0)$ & $(1,2,3)$ & $387055/369512$\\
23 & $(7,7,7);(15,5,5,1,5,1,1,0)$ & $(1,2,3)$ & $10671/9880$\\
24 & $(6,6,6);(19,4,4,0,4,0,0,2)$ & $(1,12,14)$ & $89966/88179$\\
25 & $(6,6,6);(19,4,4,0,4,0,0,2)$ & $(1,13,15)$ & $25888/24871$\\
\bottomrule
\end{tabular}
\end{center}
Consequently, both trace families forced by the 32-edge $v=25$ selection
fail when evaluated at 33 edges, and enlarging only the fixed position menus
cannot turn that safe-profile argument into a proof of $m(5)\ge34$.
\end{proposition}

The exhaustive counts and exact fractions in Proposition~\ref{prop:ceiling}
are regenerated by the optimization audit in the artifact.  The proposition
does \emph{not} exclude a proof of $m(5)\ge34$ using stronger realizability
constraints, a larger selected set, sharper counting of compatible witness
edge pairs, or a different colouring argument.  Rather, it identifies why
further tuning of the current position menus is mathematically insufficient.

\section{Reproducibility and independent verification}\label{sec:artifact}

\subsection{Artifact contents and exact outputs}

The artifact contains three primary programs under \texttt{theory/}:
\begin{itemize}[leftmargin=2em]
\item \texttt{locked\_greedy\_32.py}, implementing
      \eqref{eq:free-last}--\eqref{eq:Q};
\item \texttt{close\_v21\_v22.py}, emitting one JSONL record for every
      initial or Ramsey-step profile at $v=21,22$; and
\item \texttt{selection\_certificates.py}, regenerating every
      $v=23,24,25$ profile, clique partition, selected assignment and exact
      bound.
\end{itemize}
The $v=21,22$ artifact has 3,081 records.  Its canonical internal SHA-256
and the SHA-256 of the JSONL file are recorded in
\texttt{theory/certificates/v21\_v22\_manifest.json} and in the frozen
release checksum file.  The compact manifest
\texttt{theory/SELECTION\_MANIFEST.json} pins the six $K_4$ partitions, 32
$K_5$ partitions, every $v=23$ profile and residual set, and every
$v=23,24,25$ selected certificate row.  These latter finite sets are
regenerated from compact source rather than archived as a second large
row-by-row certificate.

The exact placement sweeps behind Proposition~\ref{prop:ceiling} and the
$v=21$ candidate-menu minimality statement, including the full labelled
counts, automorphism quotients and exact minima, are documented in
\texttt{optimization/OPTIMIZATION\_AUDIT.md} and replayed by the accompanying
optimization verifiers.

A second implementation under \texttt{solver\_alt/} does not import the
primary evaluator.  It re-derives the placement counts, enumerates
three-lock profiles from different loop variables, and generates clique
partitions as set partitions of $E(K_s)$ followed by a clique predicate.  It
replays all 3,081 $v=21,22$ rows and every $v=23,24,25$ assertion, and
compares the canonical row hashes.  Hashes are integrity checks for the
finite objects; they are not substitutes for the mathematical arguments in
Sections~\ref{sec:normal}, \ref{sec:locked}, and \ref{sec:cases}.

\subsection{Commands and reference environment}

From the artifact root, the lower-bound certificate is regenerated with
\begin{verbatim}
python3 theory/locked_greedy_32.py --self-test
python3 theory/close_v21_v22.py --output-dir theory/certificates
python3 theory/selection_certificates.py > theory/SELECTION_MANIFEST.json
python3 solver_alt/verify_v21_v22.py \
  --artifact-dir theory/certificates
python3 solver_alt/verify_lower33.py
\end{verbatim}
The position-optimality and 33-edge ceiling calculations are regenerated
from \texttt{optimization/} with
\begin{verbatim}
python3 -O general_locked.py
python3 -O compiled_locked.py
PYTHONPATH=. python3 -O verify_v25_selection_ceiling.py \
  --output v25_selection_ceiling_verification.json
PYTHONPATH=. python3 -O sweep_m32_v21_menu.py \
  --summary m32_v21_menu_summary.json \
  --records m32_v21_menu_records.jsonl
PYTHONPATH=. python3 -O verify_m32_v21_base_subsets.py \
  --summary m32_v21_base_subsets_summary.json \
  --records m32_v21_six_base_profiles.jsonl
PYTHONPATH=. python3 -O frontier_audit.py \
  --output optimization_results.json
\end{verbatim}
No third-party Python package is required.  A reference replay used CPython
3.13.1 on macOS/Darwin 25.5 (arm64), an 8-core Apple M3 and 16 GB RAM.  The
programs are single-process and exact-integer performance is sensitive to
Python version and CPU scheduling.  No minimum-memory claim is made; the
complete pipeline was verified within 16 GB.  The one-command driver prints
each command, verifies the checksum files, and reports the Python version and
total wall time.  Its pinned reports deliberately omit host-dependent paths
and timings so that byte-for-byte comparisons remain portable.

\subsection{Scope and limitations of the computer check}

The finite proof has four explicit trust boundaries.
\begin{enumerate}[leftmargin=2em]
\item Python's arbitrary-precision integers, \texttt{fractions.Fraction} and
      \texttt{math.comb} are trusted; no proof-assistant kernel checks the
      arithmetic.
\item The primary and independent implementations share only the displayed
      mathematics and pinned finite outputs, but they run on the same Python
      runtime.  Their algorithmic independence reduces coding risk without
      eliminating platform risk.
\item The profile universes are safe supersets.  Unrealizable profiles can
      make the calculation harder but cannot create a false strict
      certificate for a realizable obstruction.
\item The computational propositions verify the finite inequalities; they
      do not replace the human-checkable selection arguments that connect a
      hypothetical hypergraph to those finite families.
\end{enumerate}

\paragraph{Artifact availability.}
The code, exact certificates and verification reports supporting this article
are available from the public GitHub repository
\href{https://github.com/AyushAg1002/property-b-m5-lower33}%
{\texttt{property-b-m5-lower33}} (owner \texttt{AyushAg1002}) under release tag
\texttt{v1.0.0}.

\noindent The corresponding archival release is identified by Zenodo DOI
\href{https://doi.org/10.5281/zenodo.22070117}%
{\texttt{10.5281/zenodo.22070117}}.  Source code is licensed under the MIT
License; the manuscript source, proof notes, documentation, certificates,
manifests and verification reports are licensed under CC BY 4.0.  The
DOI-bearing release contains the exact versions of all sources, certificate
files, manifests, independent reports and optimization-audit materials cited
in this paper.

\section{Concluding remarks}

Theorem~\ref{thm:main} narrows the fifth Property-B interval to
\[
33\le m(5)\le51.
\]
The central methodological point is that a very small selected trace can
carry enough information to certify all remaining vertex orders, provided
that exact position-dependent probabilities are combined with degree and
codegree selection.  Proposition~\ref{prop:ceiling} also shows where the
present certificate families stop: a proof of $m(5)\ge34$ will require
additional structural information or a stronger probability estimate, not
merely a larger menu of locked placements for those families.  The classical
upper bound $m(5)\le51$ remains unchanged; we make no upper-bound claim.

\section*{Declarations}

\paragraph{Use of generative AI and author responsibility.}
OpenAI Codex assisted with research and manuscript preparation; the author
independently verified all proofs, computations, code, citations and wording
and accepts full responsibility for the article.

\paragraph{Conflicts of interest.}
The author declares no competing interests.

\paragraph{Funding.}
No external funding was received for this work.

\bibliographystyle{plain}
\bibliography{references}

\end{document}
